\label{cap:introducao}

A adoção cada vez maior de sistemas baseados em \gls{IA}tem alterado de maneira notável as práticas de desenvolvimento de software e os processos de tomada de decisão em vários setores essenciais. Esse crescimento é fomentado pela junção de três fatores: o aumento da disponibilidade de grandes volumes de dados, evolução das capacidades computacionais e desenvolvimento de técnicas de aprendizado de máquina. Como consequência, algoritmos de \acrshort{IA} estão cada vez mais integrados em situações onde suas decisões atuam sobre direitos fundamentais, acesso a serviços essenciais e oportunidades de vida de indivíduos e comunidades.

No entanto, o avanço tecnológico tem sido acompanhado por preocupações cada vez maiores em relação às implicações éticas dos sistemas de \acrshort{IA}. Diversos incidentes documentados evidenciam problemas sérios decorrentes da ausência de análises éticas sistemáticas durante o desenvolvimento desses sistemas. Um caso emblemático foi divulgado em 2019 por pesquisadores da Universidade de Berkeley, em colaboração com a Universidade de Chicago e a Universidade de Boston. O estudo revelou que um algoritmo amplamente utilizado por hospitais norte-americanos para alocação de cuidados de saúde, abrangendo aproximadamente 200 milhões de pacientes, apresentava um viés racial substancial. Desenvolvido pela empresa Optum, o sistema utilizava os custos históricos de saúde como um \textit{proxy} para avaliar necessidades médicas. Contudo, devido a desigualdades socioeconômicas que resultam em menores gastos médicos entre pacientes negros, o algoritmo exigia que esses indivíduos estivessem significativamente mais doentes do que pacientes brancos para receberem cuidados equivalentes. Como consequência, os cuidados recomendados para populações negras com necessidades semelhantes foram reduzidos em mais de 50\% \cite{Ziad2019}.

Problemas semelhantes ocorrem em sistemas de reconhecimento facial. Estudos conduzidos pelo \textit{National Institute of Standards and Technology} (NIST), que analisaram 189 algoritmos de 99 desenvolvedores, incluindo grandes empresas de tecnologia, revelaram que esses sistemas apresentam taxas de erro entre 10 e 100 vezes maiores na identificação de rostos de pessoas negras ou asiáticas, em comparação com rostos de pessoas brancas \cite{frvt2019}. Em particular, para mulheres negras, certos algoritmos atingem taxas de erro superiores a 35\%, enquanto, para homens brancos, esse índice é inferior a 1\%. Tais disparidades têm implicações práticas graves. Em 2020, Robert Williams foi erroneamente preso em Detroit após ser identificado por um sistema de reconhecimento facial, tornando-se o primeiro caso documentado de prisão incorreta causada por essa tecnologia nos Estados Unidos \cite{hill2020}.

Esses problemas estão diretamente relacionados à ausência de uma operacionalização sistemática de princípios éticos no desenvolvimento de sistemas de \acrshort{IA}. O viés no algoritmo da Optum poderia ter sido mitigado se princípios como justiça, não maleficência e responsabilidade tivessem sido incorporados desde as fases iniciais do desenvolvimento, incluindo a escolha crítica de variáveis \textit{proxy} e a validação com métricas desagregadas por grupos demográficos. De forma similar, os problemas em sistemas de reconhecimento facial decorrem da negligência em relação a princípios de equidade, representatividade e transparência, tanto na composição dos \textit{datasets} de treinamento quanto na comunicação das limitações conhecidas dos sistemas.

Em resposta a essas preocupações, diversos \textit{frameworks} e regulamentações internacionais têm sido propostos. O \textit{AI Act} da União Europeia \cite{EUAIAct2024} estabelece requisitos obrigatórios para sistemas de alto risco; a \textit{Recommendation on the Ethics of Artificial Intelligence}, da UNESCO \cite{UNESCO2021}, propõe princípios universais para a governança da IA; e, no Brasil, a \gls{LGPD} \cite{LGPD2018} estabelece diretrizes para o tratamento ético de dados pessoais. Esses marcos normativos reforçam a necessidade de práticas responsáveis e sinalizam uma tendência global em direção a uma regulamentação mais rigorosa.

Apesar desses avanços normativos, persiste uma lacuna significativa entre princípios éticos abstratos e sua operacionalização concreta ao longo do ciclo de vida do desenvolvimento de software. As diretrizes existentes permanecem, em sua maioria, em um nível teórico, sem oferecer soluções práticas sistematicamente integradas às diferentes fases do processo de desenvolvimento. Essa lacuna dificulta que equipes técnicas adotem práticas éticas de forma eficaz, gerando riscos de falhas éticas no produto final, custos elevados de correções posteriores e potenciais danos a usuários e comunidades.

Para fundamentar a definição precisa do problema, foi conduzida uma \gls{RSL} seguindo o protocolo de Kitchenham \& Charters \cite{barbara2007guidelines} e as diretrizes PRISMA 2020 \cite{PRISMA2020}. A busca sistemática em quatro bases de dados (ACM Digital Library, IEEE Xplore, Scopus, Web of Science) resultou na identificação inicial de 2.524 registros. Após a deduplicação e a aplicação sistemática de critérios de inclusão e exclusão em múltiplos estágios, chegou-se a um \textit{corpus} final de 75 estudos de alta qualidade (QA~$\geq$~6{,}5 em uma escala de 8{,}0), publicados entre 2021 e 2025.

A análise evidenciou oito princípios éticos predominantes, com 229 menções distribuídas entre os 75 estudos: \textit{Transparency/Explainability} (64{,}0\%), \textit{Fairness} (52{,}0\%), \textit{Accountability} (44{,}0\%), \textit{Trustworthiness} (40{,}0\%), \textit{Safety} (36{,}0\%), \textit{Privacy} (26{,}7\%), \textit{Autonomy} (16{,}0\%) e \textit{Sustainability} (4{,}0\%). Os três primeiros (\textit{Transparency/Explainability}, \textit{Fairness} e \textit{Accountability}) formam uma tríade conceitual central, identificada como a base das abordagens mais consolidadas na literatura.

A distribuição da cobertura por fases do \textit{Software Development Life Cycle} (SDLC) revelou uma concentração crítica nas fases posteriores: \textit{Implementation} (64{,}0\%) e \textit{Testing} (54{,}7\%). Em contraste, fases consideradas fundamentais para o \textit{value-sensitive design} apresentaram sub-representação severa: \textit{Requirements Engineering} (16{,}0\%), \textit{Design} (20{,}0\%), \textit{Governance} (4{,}0\%) e \textit{Maintenance} (6{,}7\%). Essa distribuição indica um padrão predominantemente reativo, em que questões éticas são tratadas principalmente durante ou após a codificação, contrastando com a abordagem proativa \textit{by design} preconizada por regulações contemporâneas.


A síntese da \acrshort{RSL} permitiu a identificação de seis lacunas críticas que caracterizam o problema investigado:

\begin{enumerate}
    \item \textbf{Sub-representação da Engenharia de Requisitos (L1):} Apenas 16{,}0\% dos estudos abordam a fase de requisitos, apesar de seu papel fundamental no \textit{value-sensitive design}. A literatura concentra-se em fases posteriores, indicando um tratamento reativo das questões éticas, em detrimento de uma abordagem preventiva;
    
    \item \textbf{Ausência de suporte ferramental integrado (L2):} Somente 6{,}7\% dos estudos propõem ferramentas integradas a ambientes de desenvolvimento (IDEs, CI/CD, MLOps). Desenvolvedores carecem de suporte prático e contextualizado para operacionalizar princípios éticos nos fluxos de trabalho cotidianos;
    
    \item \textbf{Negligência de manutenção e governança (L3):} \textit{Maintenance} (6{,}7\%), \textit{Operation} (22{,}7\%) e \textit{Governance} (4{,}0\%) estão criticamente sub-representadas, apesar da necessidade documentada de monitoramento contínuo de sistemas em produção e da existência de estruturas organizacionais de supervisão;
    
    \item \textbf{Desconexão academia--indústria (L4):} Há ausência total (0{,}0\%) de evidências de adoção industrial documentada. Faltam estudos sobre custos, benefícios, barreiras organizacionais e sustentabilidade das abordagens propostas em ambientes produtivos;
    
    \item \textbf{Trade-offs éticos não formalizados (L5):} Apenas 9{,}6\% dos estudos discutem explicitamente tensões entre princípios (por exemplo: transparência vs. privacidade; justiça vs. acurácia). A ausência de métodos estruturados para negociação e documentação de decisões éticas compromete a auditabilidade;
    
    \item \textbf{Sustentabilidade socioambiental negligenciada (L6):} Com apenas 4{,}0\% de menções, os princípios de sustentabilidade permanecem marginalizados, apesar da crescente preocupação com impactos ambientais (pegada de carbono de modelos) e sociais (deslocamento laboral, concentração de poder) de sistemas de \acrshort{IA}.
\end{enumerate}

Essas lacunas evidenciam que, apesar do crescente corpo de conhecimento sobre ética em \acrshort{IA}, exemplificado pelo aumento de 575\% no volume de publicações entre 2021 e 2025 (de 4 para 27 estudos de alta qualidade por ano), persiste uma desconexão crítica entre princípios abstratos e práticas concretas de Engenharia de Software. Profissionais carecem de orientações práticas, sistemáticas e integradas ao fluxo de trabalho para identificar, especificar e verificar requisitos éticos. Essa ausência pode resultar em sistemas que, mesmo tecnicamente avançados, falham em atender às expectativas sociais, regulatórias e éticas, gerando riscos de danos a usuários, responsabilização legal e perda de confiança pública.

Este trabalho propõe o desenvolvimento do \gls{REAI}, um modelo sistemático para a operacionalização de princípios éticos de \acrshort{IA} no contexto da Engenharia de Requisitos, com foco específico em endereçar as lacunas L1 (sub-representação da Engenharia de Requisitos), L2 (ausência de suporte ferramental integrado) e L5 (\textit{trade-offs} éticos não formalizados).

%%%%%%%%%%%%%%%%%%%%%%%%%%%%%%%%%%%%%%%%%%%%%%%%%%%%%%%%%%%%%%%%%%%%%%%%%%%%%%%%%%%%%

\section{Justificativa}

A necessidade de uma operacionalização sistemática da ética em \acrshort{IA} decorre dos desafios encontrados em equilibrar a inovação tecnológica, o direito da sociedade de beneficiar-se dos avanços científicos e a mitigação de danos potenciais associados a sistemas mal projetados. A \acrshort{RSL} conduzida nesta pesquisa fornece uma fundamentação empírica robusta para essa necessidade.

O crescimento de 575\% no volume de publicações sobre ética em \acrshort{IA} aplicada à Engenharia de Software, entre 2021 e 2025, evidencia a maturação da área e o reconhecimento crescente de sua importância. Os 75 estudos de alta qualidade analisados demonstram um esforço científico significativo na operacionalização de princípios éticos, com destaque para a tríade conceitual central: \textit{Transparency/Explainability} (64{,}0\%), \textit{Fairness} (52{,}0\%) e \textit{Accountability} (44{,}0\%). Essa concentração indica uma convergência em torno de princípios considerados fundamentais pela comunidade acadêmica e profissional.

No entanto, a análise revela uma discrepância crítica entre os aspectos teóricos e práticos. Enquanto 64{,}0\% dos estudos abordam \textit{Transparency/Explainability} como princípio, apenas 16{,}0\% propõem abordagens para a fase de \textit{Requirements Engineering} etapa em que os valores éticos deveriam ser explicitados, negociados e operacionalizados. Essa sub-representação é especialmente problemática, considerando que \textit{Requirements} constitui uma fase fundamental para o \textit{value-sensitive design}, momento em que princípios abstratos devem ser traduzidos em especificações concretas e verificáveis.

A concentração nas fases posteriores do \acrshort{SDLC}, \textit{Implementation} (64{,}0\%) e \textit{Testing} (54{,}7\%), indica um padrão predominantemente reativo. Questões éticas são tratadas principalmente durante ou após a codificação, momento em que as correções tornam-se significativamente mais custosas. Evidências consolidadas demonstram que os custos de correção aumentam exponencialmente quando problemas são identificados tardiamente \cite{Boehm2001}, reforçando a importância crítica da incorporação de considerações éticas desde as fases iniciais.

A análise da \acrshort{RSL} identificou ausência total (0{,}0\%) de evidências de adoção industrial documentada, revelando um \textit{gap} crítico entre a pesquisa acadêmica e a prática profissional. Faltam estudos sobre custos reais de implementação, benefícios mensuráveis observados, barreiras organizacionais concretas e sustentabilidade das abordagens propostas. Essa desconexão academia-indústria (L4) sugere que as soluções apresentadas podem não estar adequadamente alinhadas com as restrições e necessidades de ambientes produtivos.

Adicionalmente, apenas 6{,}7\% dos estudos propõem ferramentas integradas a ecossistemas de desenvolvimento (\textit{IDEs, CI/CD, MLOps}). Desenvolvedores necessitam de suporte prático e contextualizado, integrado aos fluxos de trabalho existentes, e não de soluções \textit{standalone} que imponham processos paralelos. Essa lacuna (L2) representa uma barreira significativa para a adoção efetiva de práticas éticas em escala industrial.

A formalização de \textit{trade-offs} éticos permanece criticamente negligenciada, com apenas 9{,}6\% dos estudos discutindo explicitamente tensões entre princípios. \textit{Trade-offs}, como transparência versus privacidade ou justiça versus acurácia, não são problemas técnicos \textit{resolvíveis} por meio de otimizações matemáticas, mas sim questões que exigem negociação participativa, explicitação de valores, documentação de decisões e mecanismos de revisão. A ausência de métodos estruturados para essa negociação (L5) compromete a auditabilidade e a responsabilização.

Em modelos de ciclo de vida de software, sejam em abordagem cascata ou ágil, a elicitação e a análise de requisitos constituem as primeiras etapas, caracterizando-se pela definição das bases do projeto, identificação de necessidades e expectativas de \textit{stakeholders}, incluindo requisitos não funcionais necessários para garantir conformidade com padrões éticos. Ao integrar princípios éticos desde a fase inicial, é possível assegurar que esses sejam considerados no design e desenvolvimento de forma estruturada e contínua. Além disso, \textit{Requirements Engineering} constitui uma fase de comunicação crítica entre desenvolvedores, usuários e diversos \textit{stakeholders}.

Propor um modelo aplicado diretamente à análise e especificação de requisitos facilita o entendimento e a aplicação prática de princípios éticos por \textit{stakeholders} com diferentes níveis de conhecimento técnico. Isso reduz o risco de falhas éticas no produto final e diminui os custos de correções posteriores, que tendem a ser significativamente elevados quando identificadas em fases avançadas do ciclo de vida ou após a implantação.

No contexto brasileiro, a justificativa adquire relevância adicional ao considerar: (i) a obrigatoriedade de conformidade com a \acrshort{LGPD}, que exige mecanismos sistemáticos para garantir o tratamento ético de dados pessoais; (ii) a necessidade de fortalecer as capacidades nacionais em desenvolvimento responsável de \acrshort{IA}, posicionando o Brasil em discussões globais sobre governança de \acrshort{IA}; (iii) a escassez de pesquisas que traduzam princípios éticos em práticas concretas de Engenharia de Software adaptadas à realidade brasileira; e (iv) a oportunidade de contribuir para debates globais com perspectivas contextualizadas às particularidades socioculturais, econômicas e regulatórias do país.


%%%%%%%%%%%%%%%%%%%%%%%%%%%%%%%%%%%%%%%%%%%%%%%%%%%%%%%%%%%%%%%%%%%%%%%%%%%%%%%%%%%%%

\section{Objetivos}
\label{sec:objetivos}

\subsection{Objetivo Geral}

O objetivo geral deste trabalho é propor e especificar o \acrshort{REAI}, um modelo de requisitos éticos para sistemas baseados em Inteligência Artificial, com foco na Engenharia de Requisitos, que endereça lacunas identificadas na literatura e mitiga riscos decorrentes da não implementação sistemática de princípios éticos.

\subsection{Objetivos Específicos}

Para alcançar o objetivo geral, esta pesquisa define os seguintes objetivos específicos:

\begin{enumerate}
    \item Conduzir uma revisão sistemática de literatura para mapear e analisar os princípios e requisitos éticos de \acrshort{IA} discutidos na literatura científica e em normativas internacionais e nacionais aplicáveis ao desenvolvimento de software;
    
    \item Identificar e caracterizar lacunas, desafios e riscos decorrentes da ausência de incorporação sistemática de princípios éticos ao longo do ciclo de vida do software, com ênfase na fase de Engenharia de Requisitos;
    
    \item Propor o modelo \acrshort{REAI}, com especificação detalhada de seus componentes (catálogo, \textit{templates}, processos, \textit{checklists}) que operacionalizam os princípios éticos identificados na literatura.
\end{enumerate}


%%%%%%%%%%%%%%%%%%%%%%%%%%%%%%%%%%%%%%%%%%%%%%%%%%%%%%%%%%%%%%%%%%%%%%%%%%%%%%%%%%%%%

\section{Resultados Esperados}

Espera-se, com este trabalho, alcançar os seguintes resultados:

\begin{enumerate}
    \item Mapeamento sistemático de processos, ferramentas, metodologias, métodos, \textit{frameworks} ou técnicas voltadas ao desenvolvimento ético de sistemas baseados em \acrshort{IA}, em diferentes fases do ciclo de vida, por meio de uma \acrshort{RSL} que analisou 75 estudos de alta qualidade publicados entre 2021 e 2025;
    
    \item Caracterização detalhada de seis lacunas críticas (L1--L6) na operacionalização de ética em \acrshort{IA}, com análise quantitativa da distribuição por fases do \textit{Software Development Life Cycle} (SDLC) e por princípios éticos, orientando agendas futuras de pesquisa;
    
    \item Especificação completa do \acrlong{REAI}, um modelo sistemático e prático para a incorporação de princípios éticos na \textit{Requirements Engineering}, que endereça, especificamente, as lacunas L1 (sub-representação da fase de requisitos, 16{,}0\%), L2 (ausência de ferramentas integradas, 6{,}7\%) e L5 (trade-offs éticos não formalizados, 9{,}6\%);
    
    \item Desenvolvimento de um conjunto de recursos práticos reutilizáveis, incluindo: catálogo de requisitos éticos, \textit{templates} de especificação, processo de elicitação, \textit{checklist} de verificação e matriz de rastreabilidade, que podem ser adaptados e estendidos por pesquisadores e profissionais de Engenharia de Software;
    
    \item Avaliação da completude, aplicabilidade e usabilidade do modelo proposto, por meio de análise qualitativa estruturada, com identificação de pontos de melhoria e direcionamentos para pesquisas futuras;
    
    \item Contribuição para o fortalecimento da comunidade brasileira no desenvolvimento responsável de \acrshort{IA}, oferecendo um recurso educacional contextualizado à realidade nacional, facilitando a conformidade com a \acrshort{LGPD} e posicionando a pesquisa brasileira em discussões globais sobre ética em \acrshort{IA}.
\end{enumerate}

Como contribuição principal, espera-se que o modelo \acrshort{REAI} forneça suporte prático e sistemático, por meio de artefatos que possam ser utilizados por engenheiros de requisitos e desenvolvedores, independentemente de experiência prévia com ética em \acrshort{IA}. Espera-se, também, que o modelo facilite a rastreabilidade e a auditoria, por meio do mapeamento explícito entre princípios éticos, requisitos especificados e decisões de design, apoiando a conformidade regulatória com a \acrshort{LGPD} e o \textit{AI Act}.

Além disso, espera-se que o modelo \acrshort{REAI} seja disponibilizado publicamente em um repositório no GitHub, sob licença aberta, incluindo documentação completa, exemplos de aplicação e materiais educacionais complementares, facilitando sua adoção e evolução colaborativa pela comunidade de Engenharia de Software.


%%%%%%%%%%%%%%%%%%%%%%%%%%%%%%%%%%%%%%%%%%%%%%%%%%%%%%%%%%%%%%%%%%%%%%%%%%%%%%%%%%%%%

\section{Método de Pesquisa}
\label{sec:metodologia}

Esta pesquisa é guiada pelo método \textit{Design Science Research} (\acrshort{DSR}) \cite{Hevner2004}. O \acrshort{DSR} é um método baseado em três conceitos fundamentais: pesquisa, design e ciência do design. A pesquisa pode ser definida como a atividade que contribui para o entendimento de um fenômeno. O design está relacionado à criação de um artefato que ainda não existe. A ciência do design é definida como o conhecimento sob a forma de construções, técnicas, métodos e modelos que permitem o mapeamento do espaço funcional para o espaço de atributos. Essa combinação de conceitos dá origem ao método \acrshort{DSR}, que pode ser descrito como uma forma de pesquisa que cria conhecimento por meio do design, da análise, da reflexão e da abstração. Portanto, o método pode ser resumido como a utilização do design como método de pesquisa, gerando conhecimento por meio do desenvolvimento de artefatos.

A utilização deste método é justificada por sua capacidade de auxiliar pesquisadores na construção e aprimoramento de artefatos, por meio de um processo contínuo de avaliação. O método é composto por um modelo de processo com cinco etapas principais e suas respectivas saídas. A Figura~\ref{fig:metodologia} apresenta as etapas e uma visão geral dos resultados desenvolvidos em cada fase. A Tabela~\ref{tab:metodologia-fases} descreve detalhadamente as características de cada etapa.

As três primeiras etapas (\textit{Consciência do Problema}, \textit{Sugestão} e \textit{Desenvolvimento}) constituem o escopo deste projeto de dissertação. As etapas subsequentes (\textit{Avaliação} e \textit{Conclusão}) serão desenvolvidas após a aprovação na banca de qualificação. Cada etapa e sua aplicação neste trabalho são apresentadas a seguir:


%----------------------------------------------------------------------------------------%
 
\begin{figure}[htbp]
  \centering
  \begin{tikzpicture}[
    scale=0.82,
    every node/.style={transform shape, font=\small, align=center},
    node distance=0.9cm,
    phase/.style={rectangle, rounded corners, minimum width=3.0cm, minimum height=0.9cm,
                  draw=black, fill=gray!10, thick},
    activity/.style={rectangle, minimum width=6.2cm, minimum height=0.55cm,
                     text width=6.0cm, draw=black!40, fill=white, font=\footnotesize},
    method/.style={rectangle, rounded corners, minimum width=6.2cm, minimum height=0.5cm,
                   text width=6.0cm, draw=black!30, fill=gray!5, font=\footnotesize},
    arrow/.style={-{Stealth[length=3mm]}, thick}
]

% FASE 1
\node[phase] (fase1) {\textbf{FASE 1}\\Consciência do Problema};
\node[activity, below=0.18cm of fase1] (ativ1) {Mapear princípios éticos em IA\\Identificar lacunas e desafios na literatura};
\node[method, below=0.15cm of ativ1] (met1) {\textit{Saída:} RSL (PRISMA 2020); Estado atual e gaps (L1–L6)};

% FASE 2
\node[phase, below=0.55cm of met1] (fase2) {\textbf{FASE 2}\\Sugestão};
\node[activity, below=0.18cm of fase2] (ativ2) {Analisar lacunas identificadas\\Especificar requisitos do artefato\\Definir escopo do modelo REAI};
\node[method, below=0.15cm of ativ2] (met2) {\textit{Saída:} Análise qualitativa; Proposta e planejamento};

% FASE 3
\node[phase, below=0.55cm of met2] (fase3) {\textbf{FASE 3}\\Desenvolvimento};
\node[activity, below=0.18cm of fase3] (ativ3) {Especificar o modelo REAI\\Desenvolver templates e checklists\\Elaborar processo de elicitação};
\node[method, below=0.15cm of ativ3] (met3) {\textit{Saída:} Modelagem conceitual; Primeira versão do modelo};

% FASE 4
\node[phase, below=0.55cm of met3] (fase4) {\textbf{FASE 4}\\Avaliação};
\node[activity, below=0.18cm of fase4] (ativ4) {Avaliar completude do modelo\\Testar aplicabilidade prática\\Verificar usabilidade dos artefatos};
\node[method, below=0.15cm of ativ4] (met4) {\textit{Saída:} Avaliação estruturada; Medidas de performance};

% FASE 5
\node[phase, below=0.55cm of met4] (fase5) {\textbf{FASE 5}\\Conclusão e Comunicação};
\node[activity, below=0.18cm of fase5] (ativ5) {Analisar resultados\\Identificar melhorias\\Documentar e publicar resultados};
\node[method, below=0.15cm of ativ5] (met5) {\textit{Saída:} Documentação final; Publicações; Defesa};

% SETAS
\draw[arrow] (met1.south) -- (fase2.north);
\draw[arrow] (met2.south) -- (fase3.north);
\draw[arrow] (met3.south) -- (fase4.north);
\draw[arrow] (met4.south) -- (fase5.north);

% RESULTADO FINAL (compacto)
\node[
    below=0.5cm of met5,
    rectangle, rounded corners,
    draw=black, fill=gray!10,
    minimum width=7cm, minimum height=0.7cm,
    text width=6.8cm, text centered
] (resultado) {
    \textbf{Resultado Final:} Modelo REAI validado\\
    \textit{Contribuição: Engenharia de Requisitos para IA Ética}
};

\end{tikzpicture}

  \caption{Etapas da Design Science Research aplicada à pesquisa}
  \label{fig:metodologia}
\end{figure}

%----------------------------------------------------------------------------------------%

\begin{table}[htb!]
  \centering
  \caption{Fases da metodologia \textit{Design Science Research} e respectivas atividades}
  \label{tab:metodologia-fases}
  \footnotesize
  \begin{tabular}{>{\bfseries}p{3cm} p{5cm} p{5cm}}
    \hline
    Fase & \textbf{Atividades} & \textbf{Métodos/Saídas} \\
    \hline

    Consciência do Problema & 
    Mapear princípios éticos em \acrshort{IA}; identificar lacunas e desafios na literatura & 
    Revisão Sistemática da Literatura (PRISMA 2020); diagnóstico do estado atual; identificação de gaps (L1 a L6) \\

    \hline

    Sugestão & 
    Analisar lacunas identificadas; especificar requisitos do artefato; definir escopo do modelo \acrshort{REAI} & 
    Análise qualitativa; proposta inicial; planejamento conceitual \\

    \hline

    Desenvolvimento & 
    Especificar o modelo \acrshort{REAI}; desenvolver templates e checklists; elaborar processo de elicitação & 
    Design Thinking; modelagem conceitual; primeira versão do modelo \\

    \hline

    Avaliação & 
    Avaliar completude do modelo; analisar aplicabilidade prática; verificar usabilidade dos artefatos & 
    Análise qualitativa; avaliação estruturada; medidas de desempenho \\

    \hline

    Conclusão e Comunicação & 
    Analisar resultados; identificar melhorias; publicar artigos; defender dissertação & 
    Documentação final; publicações científicas; apresentação e defesa \\
    
    \hline
  \end{tabular}
\end{table}


%----------------------------------------------------------------------------------------%

\subsection{Fase 1: Consciência do Problema}

Esta etapa tem como foco a conscientização do problema de pesquisa, compreendendo os atores envolvidos, os objetivos esperados, as causas do problema, seus efeitos e as contribuições potenciais ao se propor uma solução. Foi conduzida uma \gls{RSL}, seguindo o protocolo de Kitchenham \& Charters \cite{barbara2007guidelines} e as diretrizes PRISMA 2020 \cite{PRISMA2020}, com o objetivo de determinar o estado atual da ética em \acrshort{IA} e as abordagens prevalentes para sua operacionalização.

A busca sistemática em quatro bases de dados resultou na identificação inicial de 2.524 registros. Após o processo de deduplicação e a aplicação sistemática de critérios de inclusão e exclusão em múltiplos estágios, chegou-se a um \textit{corpus} final de 75 estudos de alta qualidade (QA~$\geq$~6{,}5 em uma escala de 8{,}0), publicados entre 2021 e 2025. A análise evidenciou oito princípios éticos predominantes, sendo \textit{Transparency/Explainability} (64{,}0\%), \textit{Fairness} (52{,}0\%) e \textit{Accountability} (44{,}0\%) a tríade conceitual central mais recorrente.

A distribuição de cobertura por fases do \textit{Software Development Life Cycle} (SDLC) revelou uma concentração crítica em fases posteriores, com \textit{Implementation} (64{,}0\%) e \textit{Testing} (54{,}7\%), ao passo que a fase de \textit{Requirements Engineering} apresentou sub-representação significativa (16{,}0\%). Essa análise fundamentou a identificação das seis lacunas críticas (L1 a L6) apresentadas anteriormente, justificando o foco desta pesquisa na fase de Engenharia de Requisitos.

A saída desta etapa foi a definição do estado atual das soluções práticas relacionadas à ética em \acrshort{IA} e de suas respectivas lacunas, por meio da \acrshort{RSL}, permitindo a definição e elaboração do modelo proposto. Os resultados completos encontram-se detalhados no Capítulo~\ref{cap:rsl}.

\subsection{Fase 2: Sugestão}

Subsequente à etapa anterior, esta fase corresponde ao momento em que uma nova funcionalidade é imaginada com base em uma nova configuração de elementos existentes ou inéditos. Com base nos resultados da \acrshort{RSL}, foram estabelecidos os requisitos e características iniciais do artefato a ser desenvolvido. A análise das lacunas L1 (sub-representação da fase de \textit{Requirements Engineering}, 16{,}0\%), L2 (ausência de ferramentas integradas, 6{,}7\%) e L5 (trade-offs éticos não formalizados, 9{,}6\%) orientou a decisão de propor um modelo com foco específico na Engenharia de Requisitos.

Nesta etapa, o modelo proposto foi definido como um meio para a operacionalização sistemática de requisitos éticos no contexto do desenvolvimento de sistemas baseados em \acrshort{IA}. O \acrlong{REAI} foi concebido para atender às seguintes necessidades identificadas na literatura: (i) operacionalizar princípios éticos abstratos em requisitos concretos e verificáveis; (ii) fornecer suporte sistemático desde a elicitação até a verificação de requisitos éticos; (iii) integrar-se de forma natural às práticas existentes de Engenharia de Requisitos; (iv) explicitar e documentar trade-offs entre princípios éticos potencialmente conflitantes; e (v) ser aplicável independentemente do domínio, tecnologia ou metodologia de desenvolvimento utilizada.

O resultado desta fase foi a formulação da proposta e o planejamento do modelo \acrshort{REAI}, com a especificação de cinco componentes integrados: (i) Catálogo de Requisitos Éticos, (ii) Template de \textit{User Stories} Éticas, (iii) Processo de Elicitação, (iv) Checklist de Verificação e (v) Matriz de Rastreabilidade.

\subsection{Fase 3: Desenvolvimento}

Nesta etapa, o artefato planejado é refinado e desenvolvido. Esta fase, atualmente em andamento, consiste no design e na especificação detalhada do modelo \acrshort{REAI}. É importante ressaltar que o modelo tem como objetivo facilitar a adoção de práticas éticas em sistemas baseados em \acrshort{IA} por times de desenvolvimento, com o mínimo de alterações nos fluxos já existentes de Engenharia de Software.

O modelo \acrshort{REAI} é composto por cinco componentes integrados que operacionalizam princípios éticos no contexto da Engenharia de Requisitos:

\begin{enumerate}
    \item \textbf{Catálogo de Requisitos Éticos:} mapeamento sistemático entre os oito princípios éticos identificados na \acrshort{RSL} e tipos de requisitos não funcionais correspondentes, com exemplos contextualizados por domínio de aplicação;

    \item \textbf{Template de \textit{User Stories} Éticas:} formato estruturado, adaptado de \textit{user stories}, que incorpora explicitamente princípios éticos, facilitando a especificação de requisitos em metodologias ágeis;

    \item \textbf{Processo de Elicitação:} procedimento organizado em três etapas para a identificação sistemática de requisitos éticos, considerando o contexto do projeto e os \textit{stakeholders} envolvidos;

    \item \textbf{Checklist de Verificação:} conjunto de perguntas orientadoras, organizadas por princípio ético, com o objetivo de verificar a completude e a consistência dos requisitos especificados;

    \item \textbf{Matriz de Rastreabilidade:} mecanismo de mapeamento entre requisitos éticos e os princípios que os fundamentam, facilitando auditoria e justificativa das decisões de design.
\end{enumerate}

O desenvolvimento segue uma abordagem iterativa, com refinamentos baseados na análise crítica da literatura e nas características identificadas como necessárias para o endereçamento das lacunas L1, L2 e L5. O resultado desta etapa é a primeira versão do modelo \acrshort{REAI}, apresentada no Capítulo~\ref{cap:proposta}.

\subsection{Fases 4--5: Avaliação, Conclusão e Comunicação}

As fases finais da metodologia, que serão executadas após a aprovação na banca de qualificação, estão planejadas conforme descrito a seguir.

\subsubsection{Fase 4: Avaliação} 

Nesta etapa, o artefato desenvolvido será avaliado de acordo com critérios variados, com o objetivo de gerar evidências empíricas sobre sua qualidade. O modelo \acrshort{REAI} será avaliado quanto à sua adequação teórica e aplicabilidade prática no contexto da Engenharia de Requisitos para sistemas baseados em \acrshort{IA}.

A avaliação será conduzida por meio de análise qualitativa estruturada, considerando os seguintes aspectos: (i) completude do modelo em relação à cobertura dos princípios éticos identificados na \acrshort{RSL}; (ii) aplicabilidade prática dos componentes propostos; (iii) usabilidade dos templates e checklists; (iv) efetividade na especificação de requisitos éticos verificáveis e rastreáveis; e (v) coerência interna entre os componentes do modelo.

Desvios quantitativos e qualitativos eventualmente observados serão documentados e analisados, gerando indicadores de desempenho e sugestões de melhorias. O resultado desta etapa será a avaliação do modelo proposto, com delimitação de ajustes a serem realizados em futuras iterações do desenvolvimento.

\subsubsection{Fase 5: Conclusão e Comunicação} 
Esta etapa compreende a análise final dos resultados obtidos ao longo do processo de pesquisa, identificando as contribuições teóricas e práticas do trabalho, bem como possíveis melhorias no modelo proposto e perspectivas de trabalhos futuros.

A conclusão incluirá: (i) uma síntese das contribuições alcançadas; (ii) a discussão das limitações metodológicas; (iii) recomendações para aplicação do modelo \acrshort{REAI} em contextos acadêmicos e industriais; e (iv) direcionamentos para pesquisas futuras.

A comunicação dos resultados será realizada por meio da defesa pública da dissertação, da publicação dos artigos científicos oriundos da revisão sistemática e da proposta do modelo, bem como da disponibilização do modelo em repositório público. O cronograma detalhado para execução dessas fases está apresentado no Capítulo~\ref{cap:cronograma}.


%%%%%%%%%%%%%%%%%%%%%%%%%%%%%%%%%%%%%%%%%%%%%%%%%%%%%%%%%%%%%%%%%%%%%%%%%%%%%%%%%%%%%

\section{Estrutura do Documento}

Este projeto de dissertação está organizado em seis capítulos, estruturados de acordo com as fases metodológicas apresentadas na Seção~\ref{sec:metodologia}:

\begin{itemize}
    \item \textbf{Capítulo~\ref{cap:fundamentacao} -- Referencial Teórico:} apresenta o referencial teórico relacionado à \acrshort{IA}, aos princípios de ética em \acrshort{IA}, à Engenharia de Software, à Engenharia de Requisitos e aos marcos regulatórios (LGPD, AI Act, UNESCO), contextualizando teoricamente a pesquisa;

    \item \textbf{Capítulo~\ref{cap:rsl} -- Revisão Sistemática da Literatura:} apresenta a Revisão Sistemática da Literatura conduzida neste trabalho, documentando integralmente o protocolo, a execução e os resultados, incluindo a análise de 75 estudos, a identificação de oito princípios éticos, a caracterização das lacunas críticas (L1–L6) e a síntese das evidências que fundamentam a proposta;

    \item \textbf{Capítulo~\ref{cap:proposta} -- Proposta do Modelo REAI:} apresenta a concepção, a fundamentação e a especificação detalhada do \acrlong{REAI}, incluindo a descrição de seus cinco componentes (catálogo, templates, processo, checklist e rastreabilidade), bem como a análise de sua aplicabilidade;

    \item \textbf{Capítulo~\ref{cap:cronograma} -- Cronograma Proposto:} apresenta o cronograma proposto para o desenvolvimento das fases restantes da pesquisa (avaliação, conclusão e comunicação), incluindo planejamento detalhado, marcos principais e análise de riscos;

    \item \textbf{Capítulo~\ref{cap:conclusao} -- Conclusão:} apresenta as principais conclusões e resultados do trabalho até o momento, as limitações da pesquisa, os próximos passos planejados, as melhorias potenciais e as perspectivas de continuidade.
\end{itemize}
